\section{Action Classification and Realization Theorems}
In general, we can realize each of the four prototypes, and therefore, we have a classification theorem for spaces admitting effective, Hamiltonian $\mathrm{SO}(3)$ actions, with regular momenta. To do so, we use Lemma \ref{lem:so3-sphere-hamiltonian}. In this section we 

\begin{enumerate}
    \item Prove \ref{lem:so3-sphere-hamiltonian}
    \item Give the classification theorem
\end{enumerate}

\subsection{Classification of Moment Images}

\begin{lemma}\label{lem:so3-sphere-hamiltonian}
Whenever $r>0$, the action of $\mathrm{SO}(3)$ on $S^2(r)$ given by evaluations is Hamiltonian, with moment map given by the inclusion $S^2(r) \hookrightarrow \mathbb{R}^3$ \textcite{MarsdenWeinstein1974}.
\end{lemma}
\begin{proof}
\label{MomentMapIsInclusion}
Since the coadjoint action of $\mathrm{SO}(3)$ on $\mathbb{R}^3$ is evaluation, the infinitesimal vector fields are $X^{\#}_a = \phi^{-1}(X) \times a$ where $\phi \colon \mathbb{R}^3 \to \mathfrak{so}(3)$ and $X \in \mathfrak{so}(3)$, and $a \in \mathbb{R}^3$. Recall the triple-product identity when $x \perp v$: $(a \times x) \times v = -x (a \cdot v)$. Then we note that $x \cdot ((a \times x) \times v) = -x \cdot (x (a \cdot v)) = - a \cdot v$. As we made the canonical isomorphism, we are in $\mathbb{R}^{3}$, where $\omega_{x}$ is the area form: $\omega_{x}(v,w) = x \cdot (v \times w)$. Therefore:  $\omega_{x}(X_{a}, v) = a \cdot v $. Let $\mu(x) = x$. Then: 

\begin{align}
d \langle \mu(x), a \rangle = a \cdot v &&  d \mu^{a}_{x}(v) = \omega_{x}(X_{a}, v)
\end{align}

\noindent This verifies that $\mu$ is indeed a moment map. 
\end{proof}
Now, according to Lemma \ref{lem:so3-sphere-hamiltonian}, and using Lemma \ref{lem:so3-isomorphism}, we have the following prototypes of our classification theorem, therefore, we declare using these two lemmas: 

% \begin{enumerate}
%     \item Dimension: $2$. Faithful: Yes.\\  
%     Manifold: $S^2(r)$.\\  
%     Action: $A : p \mapsto A\,p$, where $p \in S^2(r)$.\\  
%     Moment map image: $\mu(M) = S^2(r) \subset \mathbb{R}^3$ \cite{GuilleminSternberg1982}.

%     \item Dimension: $4$. Faithful: Yes.\\  
%     Manifold: $S^2(r_1) \times S^2(r_2)$, with $r_1 < r_2$.\\  
%     Action: $A : (p,q) \mapsto (A\,p,\,A\,q)$.\\  
%     Moment map image: $\{\,x \in \mathbb{R}^3 : r_2 - r_1 \leq \|x\| \leq r_1 + r_2\}$ (a closed hollow ball) \cite{MelvinParker1986}.

%     \item Dimension: $4$. Faithful: Yes.\\  
%     Manifold: $S^2(r) \times S^2(r)$.\\  
%     Action: $A : (p,q) \mapsto (A\,p,\,A\,q)$.\\  
%     Moment map image: $\{\,x \in \mathbb{R}^3 : \|x\| \leq 2r\}$ (a closed ball of radius $2r$) \cite{Blaom1996}.

%     \item Dimension: $4$. Faithful: No.\\  
%     Manifold: for example, a sphere with a trivial action.\\  
%     Action: $A : p \mapsto p$.\\  
%     Moment map image: $\mu(M) = \{0\}$ (the origin) \cite{Arnold2010}.

%     \item Dimension: $4$. Faithful: Yes.\\  
%     Manifold: $S^2(r) \times S^2(r')$.\\  
%     Action: for instance $A : (p,q) \mapsto (A\,p,\,q)$ (partial action).\\  
%     Moment map image: another sphere \cite{Silva2001}.
% \end{enumerate}

\begin{table}[h]
\small
\renewcommand{\arraystretch}{0.9}
\begin{tcolorbox}[
  enhanced,
  colback=white,
  colframe=blue!75!black,
  boxrule=0.3mm,
  width=\textwidth,
  arc=2mm,
  boxsep=1pt,
  left=2pt,
  right=2pt,
  top=2pt,
  bottom=2pt
]
\begin{tabular}{>{\centering\arraybackslash}p{0.6cm}|>{\centering\arraybackslash}p{0.6cm}|>{\centering\arraybackslash}p{2.2cm}|>{\centering\arraybackslash}p{2.5cm}|>{\centering\arraybackslash}p{8cm}}
\toprule
\textbf{$\dim$} & \textbf{F?} & \textbf{$M$} & \textbf{Action} & \textbf{$\mu(M)$} \\
\midrule

$2$ & Y & $S^2(r)$ & $A \cdot p = Ap$ & 
\begin{minipage}{8cm}
\centering
\begin{tikzpicture}[baseline, scale=0.17, tdplot_main_coords]
    \tdplotsetrotatedcoords{20}{20}{0}
    \begin{scope}[tdplot_rotated_coords]
    \draw[blue!50, thin] (0,0,0) circle (2);
    \draw[-stealth] (0,0,0) -- (2.5,0,0) node[anchor=north east,font=\tiny]{$x$};
    \draw[-stealth] (0,0,0) -- (0,2.5,0) node[anchor=north west,font=\tiny]{$y$};
    \draw[-stealth] (0,0,0) -- (0,0,2.5) node[anchor=south,font=\tiny]{$z$};
    \end{scope}
\end{tikzpicture}
\quad
{\scriptsize \textbf{Sphere} $S^2(r)$ \quad $\{x \in \mathbb{R}^3 : \|x\| = r\}$}
\end{minipage} \\ 

$4$ & Y & $S^2(r_1) \times S^2(r_2)$ & $A \cdot (p,q) = (Ap, Aq)$ & 
\begin{minipage}{8cm}
\centering
\begin{tikzpicture}[baseline, scale=0.17, tdplot_main_coords]
    \tdplotsetrotatedcoords{20}{20}{0}
    \begin{scope}[tdplot_rotated_coords]
    \draw[blue!30, fill=blue!10, fill opacity=0.3] (0,0,0) circle (2.5);
    \draw[red!30, fill=white, fill opacity=0.8] (0,0,0) circle (1);
    \draw[-stealth] (0,0,0) -- (3,0,0) node[anchor=north east,font=\tiny]{$x$};
    \draw[-stealth] (0,0,0) -- (0,3,0) node[anchor=north west,font=\tiny]{$y$};
    \draw[-stealth] (0,0,0) -- (0,0,3) node[anchor=south,font=\tiny]{$z$};
    \end{scope}
\end{tikzpicture}
\quad
{\scriptsize \textbf{Hollow ball} \quad $\{x \in \mathbb{R}^3 : r_2 - r_1 \leq \|x\| \leq r_1 + r_2\}$}
\end{minipage} \\ 

$4$ & Y & $S^2(r) \times S^2(r)$ & $A \cdot (p,q) = (Ap, Aq)$ & 
\begin{minipage}{8cm}
\centering
\begin{tikzpicture}[baseline, scale=0.17, tdplot_main_coords]
    \tdplotsetrotatedcoords{20}{20}{0}
    \begin{scope}[tdplot_rotated_coords]
    \shade[ball color=blue!40, opacity=0.4] (0,0,0) circle (2.5);
    \draw[blue!50] (0,0,0) circle (2.5);
    \draw[-stealth] (0,0,0) -- (3,0,0) node[anchor=north east,font=\tiny]{$x$};
    \draw[-stealth] (0,0,0) -- (0,3,0) node[anchor=north west,font=\tiny]{$y$};
    \draw[-stealth] (0,0,0) -- (0,0,3) node[anchor=south,font=\tiny]{$z$};
    \end{scope}
\end{tikzpicture}
\quad
{\scriptsize \textbf{Closed ball} \quad $\{x \in \mathbb{R}^3 : \|x\| \leq 2r\}$}
\end{minipage} \\ 

$2$ & N & $S^2$ & $Ap = p$ & 
\begin{minipage}{8cm}
\centering
\begin{tikzpicture}[baseline, scale=0.17, tdplot_main_coords]
    \tdplotsetrotatedcoords{20}{20}{0}
    \begin{scope}[tdplot_rotated_coords]
    \fill[black] (0,0,0) circle (0.15);
    \draw[-stealth] (0,0,0) -- (3,0,0) node[anchor=north east,font=\tiny]{$x$};
    \draw[-stealth] (0,0,0) -- (0,3,0) node[anchor=north west,font=\tiny]{$y$};
    \draw[-stealth] (0,0,0) -- (0,0,3) node[anchor=south,font=\tiny]{$z$};
    \end{scope}
\end{tikzpicture}
\quad
{\scriptsize \textbf{Origin point} \quad $\{0\}$}
\end{minipage} \\ 

$4$ & Y & $S^2(r) \times S^2$ & $A \cdot (p,q) = (Ap, q)$ & 
\begin{minipage}{8cm}
\centering
\begin{tikzpicture}[baseline, scale=0.17, tdplot_main_coords]
    \tdplotsetrotatedcoords{20}{20}{0}
    \begin{scope}[tdplot_rotated_coords]
    \draw[blue!50, thick] (0,0,0) circle (2);
    \draw[-stealth] (0,0,0) -- (3,0,0) node[anchor=north east,font=\tiny]{$x$};
    \draw[-stealth] (0,0,0) -- (0,3,0) node[anchor=north west,font=\tiny]{$y$};
    \draw[-stealth] (0,0,0) -- (0,0,3) node[anchor=south,font=\tiny]{$z$};
    \end{scope}
\end{tikzpicture}
\quad
{\scriptsize \textbf{Sphere} \quad $S^2(r)$}
\end{minipage} \\
\bottomrule
\end{tabular}
\end{tcolorbox}
\caption{\small Moment map images for different manifolds and group actions}
\end{table}

\begin{remark}
Note that there are several faithful cases.
\end{remark}

\begin{theorem}\label{thm:so3-classification}
Any compact, connected Hamiltonian $\mathrm{SO}(3)$-space with regular momenta has a moment map whose image is one of the following $\mathrm{SO}(3)$-invariant subsets of $\mathbb{R}^3$:
\begin{enumerate}
    \item A sphere $S^2(r)$ of radius $r > 0$,
    \item A closed hollow ball $\{x \in \mathbb{R}^3 : r_1 \leq \|x\| \leq r_2\}$ where $0 < r_1 < r_2$,
    \item A closed ball $\{x \in \mathbb{R}^3 : \|x\| \leq r\}$ where $r > 0$, or
    \item The origin $\{0\}$.
\end{enumerate}
Moreover, for each of these possibilities, there exists a Hamiltonian $\mathrm{SO}(3)$-space whose moment map has precisely this image.
\end{theorem}
This classification theorem, combined with our generalized Delzant correspondence (Theorem \ref{thm:delzant-so3}), provides a complete characterization of Hamiltonian $\mathrm{SO}(3)$-spaces with regular momenta, up to equivariant symplectomorphism. As shown in the previous section, the regularity condition is essential, as without it, non-equivalent spaces such as $\mathbb{CP}^2$ and $S^2(r) \times S^2(r)$ can have identical moment map images. \bigskip  

\noindent \textbf{References: } \textcite{Delzant1988, MelvinParker1986}.

\subsection{Areas of Future Inquiry}

We have that $4$-manifolds admitting effective $\mathrm{SO}(3)$-actions, but it's not easy to figure out which of them (beyond the ones we have previously mentioned) may admit symplectic forms. \textcite{MelvinParker1986} described a list of manifolds supporting Hamiltonian, effective actions: 
\begin{enumerate}
    \item $S^4$ or $\pm \mathbb{C}P^2$
    \item Connected sums of copies of $S^1 \times S^3$ and $S^1 \times P^3$
    \item $(SU(2)/H)$-bundles over $S^1$ \quad ($H$ a finite subgroup of $SU(2)$)
    \item $S^2$-bundles over surfaces
    \item Certain quotients of $S^2$-bundles over surfaces by involutions.
\end{enumerate}
Yet due to cohomological restrictions, many of these do not support a symplectic form.

\newpage

\setcounter{section}{10}